\cotefile{79}
\scanpage{113}\mspage{57}

générateurs satisfaisant

\[
\rho_0,\rho_1,\rho_\infty\quad\mid\quad \rho_0^n=\rho_1^2=\rho_\infty^3=1
\]

% slanted
\margin{\uline{NB} $n$ est exactement l'ordre de $\rho_0$ dans $G_0^+$.}

(rotations autour de $s_0$, du centre de $\{s_0,t_0\}$, et du centre de la face $\{s_0,t_0,u_0\}$), et que $G_{0\,s_0}^+$ est le sous-groupe engendré par $\rho_0$, la condition $(T_0\,3')$ est \deleted{donc} une traduction de $T_0\,3$ en termes de théorie des groupes pure. \deleted{\unsure{Dans}}

Dans les cas qui nous intéressent, où $P_0$ est régulier = en \deleted{\unsure{tant}} que polyèdre non orienté, on aura pour \deleted{\unsure{$z$|$Z$} \unsure{de}} $Z'$ commutant = à $G_0$ tout entier, non seulement à $G_0^+$, et on pourra, dans l'interprétation précédente, choisir un $g\in G_0$ -- p.ex. on pourra prendre la \deleted{symétrie}\inserted{réflexion} $\tau_0$ \deleted{\unsure{pour}} qui échange $s_0$ et $t_0$, la formule $(3\deleted{\unsure{3}}4)$ devient

\[
g\cdot\operatorname{int}(\tau_0)(gz)\ \in\ \langle\rho_0\rangle \tag*{$(34')$}
\]

\margin{sous-groupe engendré par $\rho_0$}

pour $z\in\langle\rho_0\rangle$ convenable (dépendant du choix de $g\in\operatorname{Norm}_{\deleted{\unsure{G_0^+}}}(\langle\rho_0\rangle)\unsure{)}\bigr]\unsure{]}$
% archive sheet number: 112
